% DEFINED:   (no new labels)
% RESOLVED:  cite keys: olsson2022induction, elhage2021framework,
%            wang2022wild, meng2022rome
% FORWARD:   (none)

\section*{Bibliographic Notes}

The circuit reading in this chapter follows the transformer-circuits line of
work. \textcite{elhage2021framework} sets out the framework for reading
attention transformers as compositions of QK and OV circuits, the decomposition
used throughout \Cref{sec:induction-head} to compare the pointer-lookup head to
the induction head. \textcite{olsson2022induction} identifies induction heads and
ties their emergence during pretraining to the emergence of in-context learning,
the connection \Cref{sec:icl} builds on. The pointer-dereference circuit dissected
here is a deliberately minimal cousin of those mechanisms, the same
content-addressable primitive turned on a purer memory-addressing task.

The causal interventions of \Cref{sec:scaling-the-lens} are the methodological
core of modern mechanistic interpretability. \textcite{wang2022wild} works out a
full circuit for indirect-object identification in a real language model using
activation and path patching, and makes the case that attention-pattern reading
is too weak an instrument once a computation distributes across heads and layers.
The activation-patching causal trace used on the $M = 32$ model, patching a clean
residual into a corrupted run and measuring recovery, is the technique
\textcite{meng2022rome} introduced for locating where a fact is stored in a
transformer; the same minimal-pair idea is what \Cref{fig:causal-trace} applies to
locate where the addressed bit is carried.
